\subsection{完备滤过代数中同态的提升}

设 \(k\) 为环，C 为 \(k\)-代数，\((C_n)_{n \in \mathbf{Z}}\) 为 C 的递降滤过，它与 \(k\)-代数结构相容且满足 \(C_0 = C\) (III, § 2, 第 1 节)。设 C 对由此滤过定义的拓扑是分离且完备的，于是典范映射 \(C \to \varprojlim C / C_n\) 是同胚 (同上引文, 第 6 节)。设 \(m\) 为整数 \(>0\)。以 \(\pi: C \to C / C_m\) 表示典范满射。

命题 5.—设 A 为线性拓扑的形式光滑 \(k\)-代数。每个连续 \(k\)-代数同态 \(\varphi: A \to \mathrm{C/C}_{m}\) 都有到 C 中的连续提升。

对每个整数 \(n > m\)，以 \(\pi_n: \mathrm{C/C}_n \to \mathrm{C/C}_{n-1}\) 表示典范满射。由于 C 与诸 \(\mathrm{C/C}_n\) 的逆极限恒同，给出 \(\varphi\) 到 C 的连续提升，或给出一个族 \((\varphi_n)_{n > m}\)——由连续 \(k\)-代数同态 \(\varphi_n: \mathrm{A} \to \mathrm{C/C}_n\) 组成，满足 \(\pi_n \circ \varphi_n = \varphi_{n-1}\)——两者等价。于是通过对 \(m\) 用归纳法，只需在 \(\mathrm{C}_{m+1} = 0\) 时证明陈述。此时理想 \(\mathrm{C}_m\) 的平方为零（因为 \(2m \geqslant m+1\)），由于 A 形式光滑，命题得证。

例.—设 C 为 \(k\)-代数，N 为 C 的幂零理想。该命题适用于赋予 N-进滤过的代数 C。若 A 为线性拓扑的形式光滑 \(k\)-代数，我们得到：从 A 到赋予离散拓扑的 \(k\)-代数 C/N 中的每个连续同态，都可提升为从 A 到赋予离散拓扑的 \(k\)-代数 C 中的连续同态。

